\section{The fixed seven period target}\label{sec:fixedseven}
At seven, principal units are not automatically seventh powers. The residue class of $8=1+7$ is a simple warning. The away from seven period simplification cannot be copied without an error bound. The inherited integral marked family supplies an error in $1+7^{7c}\Oloc_F$, where $c=v_7(C)\ge1$ and $F=\Q_7(\sqrt5)$. That specific error is sufficiently deep for the stated period window.

Let $X=[\phi]$ and $H=[3]$ in the mod seven local Kummer group. Exact logarithmic first digits give
\[
L(\phi)=5\overline{\sqrt5},\qquad L(3)=6,\qquad L(\sqrt5)=3,
\quad L(u)=\frac{u^{48}-1}{7}\pmod7.
\]
The two unit coordinates $X,H$ are independent, and $[\sqrt5]=4H$. For a rational seven adic unit $A$, put $q_7(A)=(A^6-1)/7$ modulo seven. The local marked period comparison has
\[
P=ZI+2XR,\qquad Z=(3q_7(A)+5)H,
\]
with $R$ from \eqref{eq:periodmatrix}.

The actual coefficient action must intertwine the toric submodule and quotient period matrices. Their valuation matrix is exactly $7cI$, so the two action matrices agree after cancellation over $\Z_7$. Their reduction commutes with $R$ and has characteristic polynomial $T^2-5$. The retained centralizer argument then gives one orientation $\epsilon\in\{1,-1\}$ and local scalar class
\begin{equation}\label{eq:kappa7}
\kappa_7=2\epsilon wX+(3q_7(A)+5)H.
\end{equation}
Merely knowing $R^2=5I$ would not identify this as the actual coefficient action. The period intertwiner is the extra geometric input.

The globally normalized natural class \eqref{eq:rigidalpha} restricts to $vwX+(r+4)H$. Comparing it with \eqref{eq:kappa7} allows one common nonzero graded scalar $g=v/(2\epsilon)$. The remaining coordinate gives
\[
2r+1=\epsilon v(3q_7(A)+5),
\]
and hence
\begin{equation}\label{eq:fermatquotient}
(2r+1)^2=v^2(3q_7(A)+5)^2\quad\text{in }\F_7.
\end{equation}
The scalar is not independently adjustable in the two unit coordinates.

For the final bounded patterns $r=6$. Patterns A and B, with $u=\pm3w$, require $q_7(A)\in\{0,6\}$ and $A^6\equiv1$ or $43$ modulo $49$. Patterns C and D, with $u=\pm w$, require $q_7(A)\in\{1,5\}$ and $A^6\equiv8$ or $36$. Each tuple still permits twelve rational unit residues $A$ modulo $49$. None is eliminated.

\subsection{Higher precision retains a valuation coordinate}
For $1\le j\le7c-1$, the deep period errors vanish in $F^\times/F^{\times7^j}$. This follows because seventh powering takes $1+7^s\Oloc_F$ bijectively onto $1+7^{s+1}\Oloc_F$ for $s\ge1$. The guaranteed window does not assert that actual errors become nontrivial at the next precision.

Use the three coordinate basis
\[
X_j=[\phi],\qquad H_j=[3],\qquad T_j=[7].
\]
The third coordinate is valuation. The actual local scalar class in the window is
\begin{equation}\label{eq:higherkappa}
\kappa_j=2\epsilon wX_j+\Lambda_jH_j+7cT_j,
\end{equation}
where
\[
\Lambda_j=7\ell_j(C_0)-3\ell_j(A)-3-5d_j,
\quad \ell_j(a)=\frac{\log(a^6)}{\log(3^6)}\pmod{7^j},
\quad d_j=\frac{\ell_j(5)}2.
\]
At level $49$, with $Q_A=(A^6-1)/7$ modulo $49$ and $q_0=(C_0^6-1)/7$ modulo seven,
\begin{equation}\label{eq:level49}
\Lambda_2=47+38Q_A+14Q_A^2-7q_0\pmod{49}.
\end{equation}
Both $C_0$ and the coordinate $7c[7]$ disappeared modulo seven. Neither may be omitted modulo forty nine.

At higher actual lattice depth, put $m=\lfloor n/2\rfloor+1$. The inherited diagonal theorem supplies an actual global upper extension $\alpha_j$ for $1\le j\le m$. In the common range $1\le j\le\min(m,7c-1)$, its local coordinates satisfy $U_j=2\epsilon wg_j$, $V_j=\Lambda_jg_j$, $W_j=7cg_j$ for one unit $g_j$. The exact valuation matrix preserves the valuation identity at every legitimate triangular precision $1\le j\le m$, even beyond the proved unit period window. Consequently
\[
v_7^{\mathrm{trunc}}(W_j)=\min(j,1+v_7(c)).
\]
At depth one, only the residual global triangular precision is supplied. Replacing the actual mod forty nine representation by an arbitrary triangular Kummer lift would erase its lower $7\beta$ entry. The local higher class exists, but its global comparison must retain the nontriangular deformation.

The archive contains exact seven adic fifth root positive controls for all admitted local residues. Their $B$ values solve $B^5=C^7-A^3$ in $\Q_7$. They are not asserted rational or ordinary integers. The higher local period criterion therefore provides real conditions while still leaving every bounded tuple locally possible.
